% !TeX root = ../tesis.tex
%%%%%%%%%%%%%%%%%%%%%%%%%%%%%%%%%%%%%%%%%%%%%%%%%%%%%%%%%%%%%%%%%%%
% §6.4 — EL ESPECTRO MJG COMPARATIVO
% Issue #25. Línea base: §5.6 (tab:espectro-mjg-diagnostico).
%
% Datos (corrida 2026-10-04; Transport-gis-zmvm-mjg/tableau/exports/
%   data/garibelt/comparativo/{kpi_comparativo,transicion_bandas,
%   distribucion_bandas}.csv):
%   C (76) 0.166 / 0.168 / 0.168 · C (141, figuras) 0.055 / 0.056 / 0.056
%   FC 0.250 (mediana 7) en los tres · T 0.696 / 0.748 / 0.822
%   B(v) 0.199 / 0.199 / 0.187 (Gini 0.8007 / 0.8012 / 0.8134)
%   DI 0.653 / 0.660 / 0.613 (muestras O-D distintas; no se compara)
%   Banda dominante por nodo: cambian 10 (MB) y 21 (Metro) nodos.
%
% Estructura (plan 2026-10-24):
%   D.1 entrada + tab:anillo-espectro ✅
%   D.2 perfil de la red en los tres casos (Unigrafica_Comparativo_
%       Valor_Normalizado, Unigrafica_Comparativo_Extendido)
%   D.3 Caso 2 (Dashboard_Espectro_MB, Dashboard_Clasificacion_MB)
%   D.4 Caso 3 (Dashboard_Espectro_Metro, Dashboard_Clasificacion_Metro)
%   D.5 clasificación por nodo (G7, Clasificacion_Banda_Dominante
%       Caso 1/MB/Metro, Comparativo_Garibelt_Cambio_Banda); definir
%       «banda dominante» aquí y, a la par, en §3.3.9
%   D.6 comparación global y límites (Dashboard_Comparativo_KPI)
%   D.7 interpretación por caso
% Decisiones del autor: DI fuera del cuadro (pies aclaran sus barras);
%   C con dominio 76 (pie aclara 141); intervalos, no adjetivos de
%   banda (#33); pies con la nomenclatura de la tesis (rótulos de las
%   figuras → issue de limpieza al terminar §6.4). Revisar #34 al
%   cerrar el capítulo.
%%%%%%%%%%%%%%%%%%%%%%%%%%%%%%%%%%%%%%%%%%%%%%%%%%%%%%%%%%%%%%%%%%%
\section{El Espectro MJG Comparativo}
\label{sec:anillo-espectro}

El Espectro MJG reúne en un solo perfil los indicadores de la red, cada uno
llevado a una escala de 0 a 1 y ubicado en uno de cuatro intervalos: 0--0.25,
0.25--0.50, 0.50--0.75 y 0.75--1.00 (§\ref{subsec:clasificacion-garibelt}).
En el Capítulo~\ref{ch:analisis-red-actual} se obtuvo el de la red actual, que
sirve aquí como línea base (§\ref{subsec:espectro-mjg}). Ninguna dimensión llega
a 0.75. La cobertura (0.166) y la intermediación (0.199) quedan en el intervalo
más bajo, la Fuerza Capilar justo en el corte de 0.25, y el tiempo de viaje
(0.696) en el de 0.50 a 0.75 (Cuadro~\ref{tab:espectro-mjg-diagnostico}). Esta
sección repite la lectura con el anillo en los Casos~2 y~3.

Con el anillo, el Espectro no se calcula por separado. Sus valores salen de los
indicadores revisados en la sección anterior (§\ref{sec:anillo-indicadores}),
de modo que solo se mueve donde se movió el indicador. La cobertura y la Fuerza
Capilar dependen de dónde están las estaciones y de con qué se enlazan. Las 94
del anillo son las mismas en los dos casos y la velocidad no forma parte de su
cálculo, así que esas dos dimensiones cambian poco. El tiempo de viaje y la
intermediación dependen de los caminos más rápidos, y esos caminos sí cambian
con la velocidad del anillo. Es ahí donde el Caso~3 se aparta del Caso~2.

Hay tres cuidados de lectura. El Factor de Desviación ($DI$) se analizó solo
para la red actual (§\ref{subsec:factor-desviacion}); si las figuras de esta
sección muestran una barra de $DI$ para los Casos~2 y~3, es porque el modelo la
calcula con otra muestra de pares origen--destino en cada caso, y por eso aquí
no se compara. La cobertura se reporta para las 76 demarcaciones de la
\zmvm{}, igual que en el Capítulo~\ref{ch:analisis-red-actual}, mientras que
las figuras usan las 141 del conjunto de datos. Por último, en el texto se da el
valor o el intervalo, no el nombre de la banda, ya que los cortes los fijó el
autor (§\ref{subsubsec:garibelt-escala}).

El Cuadro~\ref{tab:anillo-espectro} reúne los valores de las cuatro dimensiones
que se comparan, con su medida original, su valor normalizado y su intervalo
en cada caso.

\begin{table}[H]
    \centering
    \caption[Espectro MJG en los tres casos]{%
    Espectro MJG en los tres casos: medida original, valor normalizado e
    intervalo de las cuatro dimensiones que se comparan. La cobertura
    corresponde a las 76 demarcaciones de la \zmvm{}; las figuras de esta
    sección la calculan sobre las 141 del conjunto de datos (0.055 en el
    Caso~1 y 0.056 en los Casos~2 y~3). El $DI$ no se incluye porque se
    analizó solo para la red actual. El valor normalizado de la
    intermediación es $1-\mathrm{Gini}$. Corrida: 2026-10-04.}
    \label{tab:anillo-espectro}
    \small
    \setlength{\tabcolsep}{4pt}
    \begin{tabularx}{\textwidth}{|l|X|r|r|r|}
        \hline
        \textbf{Dimensión} & \textbf{Medida} & \textbf{Caso~1} &
            \textbf{Caso~2} & \textbf{Caso~3} \\
        & & \textbf{(red actual)} & \textbf{(Anillo MB)} &
            \textbf{(Anillo Metro)} \\
        \hline
        Cobertura ($C$)      & Superficie cubierta (\%) & 13.29 & 13.43 & 13.43 \\
                             & Valor normalizado        & 0.166 & 0.168 & 0.168 \\
                             & Intervalo                & 0--0.25 & 0--0.25 & 0--0.25 \\
        \hline
        Fuerza Capilar ($FC$) & Mediana de conexiones   & 7 & 7 & 7 \\
                             & Valor normalizado        & 0.250 & 0.250 & 0.250 \\
                             & Intervalo                & 0.25--0.50 & 0.25--0.50 & 0.25--0.50 \\
        \hline
        Tiempo de viaje ($T$) & Promedio (min)          & 113.91 & 108.99 & 101.92 \\
                             & Valor normalizado        & 0.696 & 0.748 & 0.822 \\
                             & Intervalo                & 0.50--0.75 & 0.50--0.75 & \textbf{0.75--1.00} \\
        \hline
        Intermediación ($B(v)$) & Índice de Gini        & 0.8007 & 0.8012 & 0.8134 \\
                             & Valor normalizado        & 0.199 & 0.199 & 0.187 \\
                             & Intervalo                & 0--0.25 & 0--0.25 & 0--0.25 \\
        \hline
    \end{tabularx}
    \fuentefigura{Elaboración propia con \texttt{VFTModel}
    \citep{cabrera2026VFTModel}; valores de §\ref{sec:anillo-indicadores}
    y del perfil comparativo de \texttt{Transport-gis-zmvm-mjg}.}
\end{table}

\subsection{El Perfil de la Red en los Tres Casos}
\label{subsec:anillo-espectro-perfil}

Como se observa en la Figura~\ref{fig:anillo-espectro-valor}, los tres casos
tienen casi el mismo perfil. Las barras de cobertura, Fuerza Capilar e
intermediación apenas se distinguen de un caso a otro. La única que se separa
con claridad es la del tiempo de viaje, que crece con el Anillo MB y crece más
con el Anillo Metro, el único caso que pasa el corte de 0.75. Con el anillo, la
red mejora en una dimensión; en las demás casi no cambia, salvo la
intermediación del Caso~3, que baja un poco.

\begin{figure}[H]
    \centering
    \includegraphics[width=\textwidth]{%
        Figures/Mapas/Tableau_Resources/Comparativo_CG/Unigrafica_Comparativo_Valor_Normalizado}
    \caption[Valor normalizado del Espectro MJG en los tres casos]{%
    Valor normalizado de cada dimensión del Espectro MJG en los tres casos,
    con líneas en los cortes de 0.25, 0.50 y 0.75. Las dimensiones
    corresponden a $C$ (accesibilidad), $B(v)$ (centralidad crítica), $FC$
    (conectividad capilar), $DI$ (eficiencia de ruta) y $T$ (fluidez global).
    La cobertura se calcula sobre las 141 demarcaciones del conjunto de datos
    (en el texto, 76), y la barra de $DI$ de los Casos~2 y~3 no se compara
    porque cada caso usa otra muestra de pares origen--destino.
    Versión a página completa en la
    Figura~\ref{fig:anx-g-espectro-valor}. Corrida: 2026-10-04.}
    \label{fig:anillo-espectro-valor}
    \fuentefigura{Elaboración propia con \textbf{Tableau};
    datos: \texttt{VFTModel}~\citep{cabrera2026VFTModel}.}
\end{figure}

La Figura~\ref{fig:anillo-espectro-cambio} muestra la diferencia de cada caso
frente a la red actual. El cambio más grande del perfil es el del tiempo de
viaje, que sube 0.052 con el Anillo MB y 0.126 con el Anillo Metro. La
intermediación es la única dimensión que retrocede: baja 0.013 en el Caso~3,
porque se concentra en las estaciones del anillo
(§\ref{subsec:anillo-centralidad}). La cobertura gana una milésima en los dos
casos y la Fuerza Capilar no cambia.

\begin{figure}[H]
    \centering
    \includegraphics[width=\textwidth]{%
        Figures/Mapas/Tableau_Resources/Comparativo_CG/Unigrafica_Comparativo_Extendido}
    \caption[Cambio del Espectro MJG frente a la red actual]{%
    Cambio del valor normalizado de cada dimensión en los Casos~2 y~3 frente
    al Caso~1; los valores positivos indican mejora. La cobertura se calcula
    sobre las 141 demarcaciones del conjunto de datos, y el cambio de $DI$ no
    se interpreta porque cada caso usa otra muestra de pares
    origen--destino. Versión a página completa en la
    Figura~\ref{fig:anx-g-espectro-cambio}. Corrida: 2026-10-04.}
    \label{fig:anillo-espectro-cambio}
    \fuentefigura{Elaboración propia con \textbf{Tableau};
    datos: \texttt{VFTModel}~\citep{cabrera2026VFTModel}.}
\end{figure}

Leído con la pregunta urbana de cada dimensión
(§\ref{subsubsec:garibelt-preguntas}), el resultado queda así. ¿Cuánto del
territorio está a una distancia caminable de la red? Casi lo mismo que hoy,
porque en la \cdmx{} el Periférico pasa por zonas que ya estaban cubiertas
(§\ref{subsec:anillo-cobertura}). ¿Qué tan variadas son las opciones de
conexión en cada nodo? Igual que antes. ¿Cuánto dura un viaje típico de un
extremo a otro? Menos, y bastante menos con el Anillo Metro. ¿Depende la red de
unos pocos nodos? Sí, en los tres casos, y con el Anillo Metro un poco más,
aunque esos nodos ya no están solo en el centro.

\subsubsection{Caso~2: Anillo MB}

En el Caso~2 el perfil es el de la red actual con una sola diferencia, como se
observa en la Figura~\ref{fig:anillo-espectro-mb}. El tiempo de viaje pasa de
0.696 a 0.748 y queda a 0.0025 del corte de 0.75. A la velocidad del Mexibús,
el anillo acorta el viaje promedio casi cinco minutos, pero no lo suficiente
para cambiar de intervalo. Las otras tres dimensiones se quedan donde estaban:
la cobertura sube de 0.166 a 0.168, y la Fuerza Capilar y la intermediación
siguen en 0.250 y 0.199.

\begin{figure}[H]
    \centering
    \includegraphics[width=\textwidth]{%
        Figures/Mapas/Tableau_Resources/RedConAnillo/Escenario_MB/Dashboard_Espectro_MB}
    \caption[Espectro MJG en el Caso~2]{%
    Espectro MJG en el Caso~2 (Anillo MB): valor normalizado de cada
    dimensión, con su medida original anotada sobre la barra. La cobertura
    (4.5\,\%) se calcula sobre las 141 demarcaciones del conjunto de datos;
    la barra de $DI$ no se compara entre casos, y los 85~min que se anotan en
    $T$ son el tiempo de referencia del modelo. Versión a página completa en la
    Figura~\ref{fig:anx-g-espectro-mb}. Corrida: 2026-10-04.}
    \label{fig:anillo-espectro-mb}
    \fuentefigura{Elaboración propia con \textbf{Tableau};
    datos: \texttt{VFTModel}~\citep{cabrera2026VFTModel}.}
\end{figure}

La Figura~\ref{fig:anillo-espectro-clasificacion-mb} junta ese perfil con las
diez estaciones de mayor intermediación por nombre, las mismas de la
Figura~\ref{fig:anillo-bv-dashboard-mb}. Las dos vistas dicen lo mismo desde
lados distintos. El perfil casi no cambia porque los nodos de los que depende
la red siguen siendo los del centro, de Tacubaya a La Raza. El Anillo MB hace
más corto el viaje, pero no toca la estructura que sostiene a la red.

\begin{figure}[H]
    \centering
    \includegraphics[width=\textwidth]{%
        Figures/Mapas/Tableau_Resources/RedConAnillo/Escenario_MB/Dashboard_Clasificacion_MB}
    \caption[Clasificación Garibelt y orden de intermediación en el Caso~2]{%
    Clasificación Garibelt en el Caso~2 (Anillo MB): perfil del Espectro MJG,
    diez estaciones con mayor suma de $B(v)$ por nombre (valor más alto de
    cada sistema) y medidas originales de cuatro dimensiones. Mismas notas
    sobre la cobertura y el $DI$ que en la
    Figura~\ref{fig:anillo-espectro-mb}. Versión a página completa en la
    Figura~\ref{fig:anx-g-espectro-clasificacion-mb}. Corrida: 2026-10-04.}
    \label{fig:anillo-espectro-clasificacion-mb}
    \fuentefigura{Elaboración propia con \textbf{Tableau};
    datos: \texttt{VFTModel}~\citep{cabrera2026VFTModel}.}
\end{figure}

\subsubsection{Caso~3: Anillo Metro}

En el Caso~3 el perfil cambia en dos dimensiones
(Figura~\ref{fig:anillo-espectro-metro}). El tiempo de viaje sube a 0.822 y
pasa al intervalo de 0.75 a 1.00, el único caso en que una dimensión del
Espectro llega ahí. La intermediación baja de 0.199 a 0.187 y sigue en el
intervalo más bajo. La cobertura y la Fuerza Capilar quedan igual que en el
Caso~2.

\begin{figure}[H]
    \centering
    \includegraphics[width=\textwidth]{%
        Figures/Mapas/Tableau_Resources/RedConAnillo/Escenario_Metro/Dashboard_Espectro_Metro}
    \caption[Espectro MJG en el Caso~3]{%
    Espectro MJG en el Caso~3 (Anillo Metro): valor normalizado de cada
    dimensión, con su medida original anotada sobre la barra. Mismas notas
    sobre la cobertura, el $DI$ y el tiempo de referencia del modelo que en
    la Figura~\ref{fig:anillo-espectro-mb}. Versión a página completa en la
    Figura~\ref{fig:anx-g-espectro-metro}. Corrida: 2026-10-04.}
    \label{fig:anillo-espectro-metro}
    \fuentefigura{Elaboración propia con \textbf{Tableau};
    datos: \texttt{VFTModel}~\citep{cabrera2026VFTModel}.}
\end{figure}

Las dos dimensiones que cambian lo hacen en sentidos opuestos, y la
Figura~\ref{fig:anillo-espectro-clasificacion-metro} ayuda a ver por qué. Junto
al perfil aparecen las diez estaciones con más intermediación por nombre.
Tacubaya sigue en el primer lugar, pero Periférico Oriente/Tláhuac e
Iztapalapa, las dos del anillo, ocupan el segundo y el tercero. El mismo anillo
que acorta los viajes junta en sus estaciones los caminos que antes cruzaban el
centro, y eso es lo que baja el valor de la intermediación. Con el Anillo
Metro, la red tarda menos, pero depende un poco más de unos cuantos nodos.

\begin{figure}[H]
    \centering
    \includegraphics[width=\textwidth]{%
        Figures/Mapas/Tableau_Resources/RedConAnillo/Escenario_Metro/Dashboard_Clasificacion_Metro}
    \caption[Clasificación Garibelt y orden de intermediación en el Caso~3]{%
    Clasificación Garibelt en el Caso~3 (Anillo Metro): perfil del Espectro
    MJG, diez estaciones con mayor suma de $B(v)$ por nombre (valor más alto
    de cada sistema) y medidas originales de cuatro dimensiones. Mismas notas
    que en la Figura~\ref{fig:anillo-espectro-mb}. Versión a página completa en la
    Figura~\ref{fig:anx-g-espectro-clasificacion-metro}. Corrida: 2026-10-04.}
    \label{fig:anillo-espectro-clasificacion-metro}
    \fuentefigura{Elaboración propia con \textbf{Tableau};
    datos: \texttt{VFTModel}~\citep{cabrera2026VFTModel}.}
\end{figure}

\subsection{La Clasificación por Nodo}
\label{subsec:anillo-espectro-nodo}

El Espectro también clasifica cada nodo por separado. A cada uno se le asigna
una banda dominante, que es el intervalo del menor de dos valores: su Fuerza
Capilar y su intermediación, normalizadas con el mínimo y el máximo de cada
caso (§\ref{subsubsec:garibelt-procedimiento}). La
Figura~\ref{fig:anillo-banda-caso1} muestra el resultado en la red actual. El
99.4\,\% de los nodos queda en el intervalo de 0 a 0.25, 68 están en el de 0.25
a 0.50 y solo dos, en Zapata y Mixcoac, llegan al de 0.50 a 0.75. Casi siempre
es la intermediación la que limita, porque se mide contra la de Tacubaya y muy
pocos nodos se le acercan. Los que salen del intervalo más bajo dibujan el
corredor central, de Garibaldi a Barranca del Muerto.

\begin{figure}[H]
    \centering
    \includegraphics[width=\textwidth]{%
        Figures/Mapas/Qgis_Resources/RedActual/Clasificacion_Garibelt/Clasificacion_Banda_Dominante}
    \caption[Banda dominante por nodo en el Caso~1]{%
    Banda dominante por nodo en el Caso~1 (red actual). Los colores
    corresponden a los intervalos 0--0.25, 0.25--0.50 y 0.50--0.75. La nota
    del mapa da 78 nodos en el intervalo de 0.25 a 0.50; los datos de la
    corrida dan 68, que es la cifra del texto. Versión a página completa en la
    Figura~\ref{fig:anx-g-banda-caso1}. Corrida: 2026-10-04.}
    \label{fig:anillo-banda-caso1}
    \fuentefigura{Elaboración propia con \textbf{QGIS};
    datos: \texttt{Apímetro} y \texttt{VFTModel}~\citep{cabrera2025apimetro, cabrera2026VFTModel}.}
\end{figure}

Con el anillo la distribución casi no cambia, como se ve en las
Figuras~\ref{fig:anillo-banda-mb} y~\ref{fig:anillo-banda-metro}. En el Caso~2
siguen en el intervalo más bajo 11,138 de los 11,209 nodos (99.4\,\%), y en el
Caso~3, 11,125 (99.3\,\%). De las 94 estaciones del anillo, tres llegan al
intervalo de 0.25 a 0.50 con el Anillo MB y 13 con el Anillo Metro. Las demás
se quedan abajo por sus pocas conexiones
(§\ref{subsec:anillo-fuerza-capilar}), incluso las que en el Caso~3 tienen
mucha intermediación.

\begin{figure}[H]
    \centering
    \includegraphics[width=\textwidth]{%
        Figures/Mapas/Qgis_Resources/RedConAnillo/Escenario_MB/Clasificacion_Garibelt/Clasificacion_Banda_Dominante_MB}
    \caption[Banda dominante por nodo en el Caso~2]{%
    Banda dominante por nodo en el Caso~2 (Anillo MB), con los mismos
    intervalos de la Figura~\ref{fig:anillo-banda-caso1}. Versión a página
    completa en la Figura~\ref{fig:anx-g-banda-mb}. Corrida: 2026-10-04.}
    \label{fig:anillo-banda-mb}
    \fuentefigura{Elaboración propia con \textbf{QGIS};
    datos: \texttt{Apímetro} y \texttt{VFTModel}~\citep{cabrera2025apimetro, cabrera2026VFTModel}.}
\end{figure}

\begin{figure}[H]
    \centering
    \includegraphics[width=\textwidth]{%
        Figures/Mapas/Qgis_Resources/RedConAnillo/Escenario_Metro/Clasificacion_Garibelt/Clasificacion_Banda_Dominante_Metro}
    \caption[Banda dominante por nodo en el Caso~3]{%
    Banda dominante por nodo en el Caso~3 (Anillo Metro), con los mismos
    intervalos de la Figura~\ref{fig:anillo-banda-caso1}. Versión a página
    completa en la Figura~\ref{fig:anx-g-banda-metro}. Corrida: 2026-10-04.}
    \label{fig:anillo-banda-metro}
    \fuentefigura{Elaboración propia con \textbf{QGIS};
    datos: \texttt{Apímetro} y \texttt{VFTModel}~\citep{cabrera2025apimetro, cabrera2026VFTModel}.}
\end{figure}

La Figura~\ref{fig:anillo-banda-transicion} cuenta los nodos existentes que
cambian de intervalo. Son 10 en el Caso~2, cuatro que suben y seis que bajan, y
21 en el Caso~3, 11 que suben y 10 que bajan. Todos se mueven entre los dos
intervalos más bajos, y en todos el cambio viene de la intermediación. Hay 39
nodos cuya Fuerza Capilar sube de intervalo en los dos casos, pero ninguno
cambia de banda dominante por esa razón, porque en ellos la intermediación
sigue siendo el valor más bajo.

\begin{figure}[H]
    \centering
    \includegraphics[width=\textwidth]{%
        Figures/Cap6/TiempoViajePromedio/Grafica_Espectro_Transicion_Banda}
    \caption[Cambio de banda dominante por nodo en los Casos~2 y~3]{%
    Nodos por intervalo de banda dominante en el Caso~1 (filas) y en los
    Casos~2 y~3 (columnas). Las celdas de color intenso son los nodos
    existentes que cambian de intervalo; la última fila corresponde a las
    94 estaciones del anillo, que no existen en el Caso~1.
    Corrida: 2026-10-04.}
    \label{fig:anillo-banda-transicion}
    \fuentefigura{Elaboración propia con \texttt{Transport-gis-zmvm-mjg};
    datos: \texttt{VFTModel}~\citep{cabrera2026VFTModel}.}
\end{figure}

La Figura~\ref{fig:anillo-banda-cambio} muestra dónde ocurren esos cambios. Con
el Anillo Metro el patrón es claro. Diez de los 11 nodos que suben están sobre
el Periférico o a menos de 1.5~km de él, como Periférico Oriente, Constitución
de 1917, Barranca del Muerto o Luis Cabrera. Nueve de los diez que bajan están
en el interior, sobre todo en la Calzada de Tlalpan y el Tren Ligero, entre
Xola y La Virgen, además de Bellas Artes y La Raza. Con el Anillo MB el efecto
es menor: cuatro nodos suben en el poniente y el oriente del Periférico y seis
bajan. Conviene una precaución. La intermediación de cada nodo se mide contra
el máximo de su caso, y ese máximo baja de 0.223 a 0.157 con el Anillo Metro;
por eso tres nodos que pierden intermediación, entre ellos UAM-I y Chabacano,
suben de intervalo. En esos tres la mejora es solo relativa.

\begin{figure}[H]
    \centering
    \includegraphics[width=\textwidth]{%
        Figures/Mapas/Qgis_Resources/Comparativo_CG/Comparativo_Garibelt_Cambio_Banda}
    \caption[Nodos que cambian de banda dominante con el anillo]{%
    Nodos existentes que cambian de banda dominante en el Caso~2 (panel~A,
    10 nodos) y en el Caso~3 (panel~B, 21 nodos) frente al Caso~1. Los
    triángulos hacia arriba pasan del intervalo de 0 a 0.25 al de 0.25 a
    0.50, y los triángulos hacia abajo, al revés. Versión a página completa en la
    Figura~\ref{fig:anx-g-banda-cambio}. Corrida: 2026-10-04.}
    \label{fig:anillo-banda-cambio}
    \fuentefigura{Elaboración propia con \textbf{QGIS};
    datos: \texttt{Apímetro} y \texttt{VFTModel}~\citep{cabrera2025apimetro, cabrera2026VFTModel}.}
\end{figure}

\subsection{Comparación Global y Límites}
\label{subsec:anillo-espectro-global}

La Figura~\ref{fig:anillo-espectro-kpi} reúne en una sola vista el perfil de los
tres casos y su cambio frente a la red actual. Leído en conjunto, el Espectro
dice que el anillo resuelve una parte del diagnóstico y deja otra como estaba.
Mejora la fluidez, que ya era la dimensión con el valor más alto, y no toca las
dos que la línea base mostraba como las más bajas, la cobertura y la
intermediación. Esas dos dependen de llegar a territorio nuevo y de repartir
los caminos entre más nodos, y un solo corredor sobre el Periférico no hace
ninguna de las dos cosas.

\begin{figure}[H]
    \centering
    \includegraphics[width=\textwidth]{%
        Figures/Mapas/Tableau_Resources/Comparativo_CG/Dashboard_Comparativo_KPI}
    \caption[\termino{Dashboard} comparativo del Espectro MJG]{%
    Vista conjunta del Espectro MJG en los tres casos: valor normalizado de
    cada dimensión (arriba) y cambio de los Casos~2 y~3 frente al Caso~1
    (abajo). La cobertura se calcula sobre las 141 demarcaciones del
    conjunto de datos, y la barra de $DI$ no se compara entre casos.
    Versión a página completa en la
    Figura~\ref{fig:anx-g-espectro-kpi}. Corrida: 2026-10-04.}
    \label{fig:anillo-espectro-kpi}
    \fuentefigura{Elaboración propia con \textbf{Tableau};
    datos: \texttt{VFTModel}~\citep{cabrera2026VFTModel}.}
\end{figure}

Como en las secciones anteriores, la comparación tiene límites que conviene
tener presentes:

\begin{itemize}
    \item \textbf{Mismos supuestos de los indicadores.} El Espectro hereda los
    límites de $T$ y de $B(v)$: no pondera por demanda, aplica la misma
    penalización por transbordo en los tres casos y supone una red sin
    contingencias.
    \item \textbf{Cortes fijados por el autor.} Los cortes de 0.25, 0.50 y
    0.75 y el rango de 85 a 180~minutos de $T$ son supuestos. El Caso~2 está
    tan cerca del corte de 0.75 que bastaría con llevar la cota de 180 a
    181~minutos para que también lo pasara.
    \item \textbf{Clasificación por nodo relativa.} La banda dominante se mide
    contra el máximo de cada caso, así que un nodo puede cambiar de intervalo
    sin que cambie su intermediación.
\end{itemize}

\subsection{Interpretación por Caso}
\label{subsec:anillo-espectro-interpretacion}

\textbf{Caso~2 (Anillo MB).} El Anillo MB deja el Espectro casi igual. Solo
mejora el tiempo de viaje, que se queda a 0.0025 del intervalo más alto. La
cobertura, la Fuerza Capilar y la intermediación conservan sus valores, y
apenas diez nodos cambian de banda dominante. Para el diagnóstico, la red con
el Anillo MB es la red actual con viajes algo más cortos: llega al mismo
territorio y depende de los mismos nodos del centro.

\textbf{Caso~3 (Anillo Metro).} El Anillo Metro es el único caso que mueve el
Espectro. El tiempo de viaje pasa al intervalo de 0.75 a 1.00, la primera
dimensión de la red que llega ahí. A cambio, la intermediación baja de 0.199 a
0.187, porque la dependencia de pocos nodos pasa del centro al Periférico. En
la clasificación por nodo, los que suben están sobre el anillo y los que bajan,
en la Calzada de Tlalpan y el centro. La red queda más rápida y con su peso
corrido hacia el sur y el oriente, pero no menos concentrada.
